\chapter{راهنمای اجرای پروژه}

این پیوست خلاصه‌ای از روش اجرای خط لوله پروژه را ارائه می‌کند. دستورهای زیر بر اساس ساختار مخزن \lr{\texttt{CE-FINAL-PROJECT}} نوشته شده‌اند.

\section{ساخت ماژول‌های جاوا}
\begin{latin}
\begin{verbatim}
mvn package -DskipTests
\end{verbatim}
\end{latin}

\section{اجرای اجزای اصلی}
ابتدا \lr{Kafka} و سرویس‌های وابسته اجرا می‌شوند. سپس دریافت‌کننده رخدادهای \lr{JFR}، دریافت‌کننده رخدادهای هسته، تجمیع‌کننده و تحلیل‌گر اجرا می‌شوند. اسکریپت \lr{\texttt{scripts/run-profiler-pipelines.sh}} برای هماهنگ کردن این مراحل در مخزن پروژه قرار دارد.

\section{اجرای ارزیابی}
برای اجرای ماژول ارزیابی:
\begin{latin}
\begin{verbatim}
mvn -pl profiler-evaluation -am package -DskipTests
java -cp profiler-evaluation/target/classes aut.ce.evaluation.App
\end{verbatim}
\end{latin}

برنامه در شروع اجرا شناسه فرایند را چاپ می‌کند. برای بررسی رخدادهای هسته، فرستنده \lr{eBPF} با همان شناسه فرایند ساخته و اجرا می‌شود:
\begin{latin}
\begin{verbatim}
cd kernel-event-sender
make clean
make TARGET_PID=12345
sudo ./build/loader
\end{verbatim}
\end{latin}

\section{خروجی‌ها}
خروجی‌های ارزیابی در قالب فایل‌های زیر تولید می‌شوند:
\begin{latin}
\begin{verbatim}
experiment_results.csv
experiment_metrics.json
evaluation_report.md
\end{verbatim}
\end{latin}

گزارش نهایی شامل وضعیت موفقیت تشخیص در هر سناریو و تأخیر تشخیص است.
